\documentclass[11pt]{article}

% Change "review" to "final" to generate the final (sometimes called camera-ready) version.
% Change to "preprint" to generate a non-anonymous version with page numbers.
\usepackage[review]{acl}

% Standard package includes
\usepackage{times}
\usepackage{latexsym}

% For proper rendering and hyphenation of words containing Latin characters (including in bib files)
\usepackage[T1]{fontenc}
% For Vietnamese characters
% \usepackage[T5]{fontenc}
% See https://www.latex-project.org/help/documentation/encguide.pdf for other character sets

% This assumes your files are encoded as UTF8
\usepackage[utf8]{inputenc}

% This is not strictly necessary, and may be commented out,
% but it will improve the layout of the manuscript,
% and will typically save some space.
\usepackage{microtype}

% This is also not strictly necessary, and may be commented out.
% However, it will improve the aesthetics of text in
% the typewriter font.
\usepackage{inconsolata}

%Including images in your LaTeX document requires adding
%additional package(s)
\usepackage{graphicx}

% If the title and author information does not fit in the area allocated, uncomment the following
%
%\setlength\titlebox{<dim>}
%
% and set <dim> to something 5cm or larger.

\title{Kelvi: a “pedagogical” morphological parser to support Tamil literacy}

% Author information can be set in various styles:
% For several authors from the same institution:
% \author{Author 1 \and ... \and Author n \\
%         Address line \\ ... \\ Address line}
% if the names do not fit well on one line use
%         Author 1 \\ {\bf Author 2} \\ ... \\ {\bf Author n} \\
% For authors from different institutions:
% \author{Author 1 \\ Address line \\  ... \\ Address line
%         \And  ... \And
%         Author n \\ Address line \\ ... \\ Address line}
% To start a separate ``row'' of authors use \AND, as in
% \author{Author 1 \\ Address line \\  ... \\ Address line
%         \AND
%         Author 2 \\ Address line \\ ... \\ Address line \And
%         Author 3 \\ Address line \\ ... \\ Address line}

\author{
 \textbf{Shankhalika Srikanth\textsuperscript{1}},
 \textbf{Sabrina Yu\textsuperscript{1}},
 \textbf{Sophia Chan\textsuperscript{1}},
 \textbf{Madeline Solis de Ovando \textsuperscript{1}},
 \textbf{Fizza Ahmed \textsuperscript{1}}
 \\
 \textsuperscript{1}Independent Researcher,
 \\
 \small{
   \textbf{Correspondence:} \href{mailto:shankhalika.srikanth@alumni.utoronto.ca}{shankhalika.srikanth@alumni.utoronto.ca}
 }
}
%\author{
%  \textbf{First Author\textsuperscript{1}},
%  \textbf{Second Author\textsuperscript{1,2}},
%  \textbf{Third T. Author\textsuperscript{1}},
%  \textbf{Fourth Author\textsuperscript{1}},
%\\
%  \textbf{Fifth Author\textsuperscript{1,2}},
%  \textbf{Sixth Author\textsuperscript{1}},
%  \textbf{Seventh Author\textsuperscript{1}},
%  \textbf{Eighth Author \textsuperscript{1,2,3,4}},
%\\
%  \textbf{Ninth Author\textsuperscript{1}},
%  \textbf{Tenth Author\textsuperscript{1}},
%  \textbf{Eleventh E. Author\textsuperscript{1,2,3,4,5}},
%  \textbf{Twelfth Author\textsuperscript{1}},
%\\
%  \textbf{Thirteenth Author\textsuperscript{3}},
%  \textbf{Fourteenth F. Author\textsuperscript{2,4}},
%  \textbf{Fifteenth Author\textsuperscript{1}},
%  \textbf{Sixteenth Author\textsuperscript{1}},
%\\
%  \textbf{Seventeenth S. Author\textsuperscript{4,5}},
%  \textbf{Eighteenth Author\textsuperscript{3,4}},
%  \textbf{Nineteenth N. Author\textsuperscript{2,5}},
%  \textbf{Twentieth Author\textsuperscript{1}}
%\\
%\\
%  \textsuperscript{1}Affiliation 1,
%  \textsuperscript{2}Affiliation 2,
%  \textsuperscript{3}Affiliation 3,
%  \textsuperscript{4}Affiliation 4,
%  \textsuperscript{5}Affiliation 5
%\\
%  \small{
%    \textbf{Correspondence:} \href{mailto:email@domain}{email@domain}
%  }
%}

\begin{document}
\maketitle
\begin{abstract}
In this paper we discuss the development of kelvi.ca, an open source web-based dictionary and morphological parser designed to aid Tamil learners in developing their literacy skills. 
Tamil is an agglutinative language and heavily suffixal. Existing Tamil dictionaries only carry stems, not conjugated or inflected forms, and for a beginner learner of the language, isolating the stem in an unfamiliar word can be very challenging. 
Kelvi provides 1) the stem of any input word alongside its definition, and 2) non-technical descriptions of any suffixes that are part of this input, so that learners will gradually start to recognize these suffixes and be able to understand and produce new Tamil words themselves. 
We also hope to show that Kelvi can be adapted for other languages and has the potential to be a useful pedagogical aid for learner literacy development, especially for agglutinative and/or polysynthetic languages which tend to be otherwise underserved in the mainstream.
\end{abstract}

\section{Introduction}

Tamil is a Dravidian language spoken primarily in south Asia, with significant global diaspora populations. 
Overall there number around 80 million Tamil speakers, with nearly 240,000 speakers in Canada alone (\citealp{government_of_canada_language_2025}, worlddata.info). 
The diaspora population in Canada largely consists of heritage speakers - that is to say, people who grow up hearing their family speak Tamil, but aren’t necessarily entirely fluent themselves \citep{polinsky_heritage_2007}. 
Two of the main complicating features of Tamil literacy are the language’s heavily diglossic nature and its agglutinative properties. 
Tamil is heavily suffixal, with most grammatical information expressed either as nominal or verbal suffixes. 
See (1), where a complete utterance is expressed as one word. 
% insert gloss of Tmail word
In its written Tamil the higher, formal register is used, where every suffix is phonologically expressed in its entirety. 
In the colloquial spoken register, however, Tamil words undergo massive phonological reduction, 
so that many sounds that are fully expressed in the written form of a word aren’t actually pronounced (\citealp{annamalai_modern_2019}; \citealp{renganathan_heritage_2015}). 
This means that heritage speakers often struggle to recognize the written version of words they have only ever heard in the colloquial spoken register, 
since the written form “sounds” different from what they’re used to \citep{renganathan_heritage_2015}. 
Conjugated words don’t appear in standard dictionaries, 
only their lemmas do (consider how “go” can be found in an English dictionary but “went” wouldn’t be.) 
Learners then face the “dictionary problem”; to summarize \citet{littell_waldayu_2017}, 
dictionaries are intended to be a tool for language learners, but paradoxically require knowledge of the language to use; 
in this case, knowledge of what the “dictionary form” or lemma of an unfamiliar word is. 
The tool we built, kelvi.ca, is a dictionary and morphological parser that helps the learner who comes 
across an intimidatingly long Tamil word see how the word is broken down, and what each “part” (morpheme) means 
in layman terms, so that they come away with a fuller understanding of the word while hopefully also starting to 
recognize Tamil morphological patterns themselves.

\section{Literature Review}
The one specific instance known to these authors of a dictionary specifically designed for a heritage/minoritized 
language that also provides pedagogical information about morphology is the ongoing work on the 
digital Wendat Dictionary \citep{lukaniec_modern_2023}. 
More broadly speaking, there are several open-source technologies that have been developed to 
support language revitalization and maintenance efforts \citep{brinklow_indigenous_nodate}, 
such as the dictionary-builder Mother Tongues \citep{littell_waldayu_2017} and the 
morphological generator Gramble \citep{littell_gramble_2024}. 
For Tamil specifically there exist technical and/or academic tools such as a lemmatizer \citep{qi2020stanza}, 
a morphological analyser \citep{sarveswaran_thamizhimorph_2021},  and a stemmer \citep{rajalingam_rdamodharantamil-stemmer_2025}, as well as some 
pedagogical resources like a dictionary of Tamil verbs \citep{schiffman_english_2009} and a few 
isolated courses/course materials (\citealp{annamalai_colloquial_2002}, \citealp{prasad_tamil_2025}, \citealp{cheran_learn_2001}); 
but to the authors’ knowledge there is currently no freely available online resource that employs Tamil NLP, 
specifically morphological analysis, for the purpose of pedagogy.

\section{Methods}
In accordance with principles behind collaborative community based research \citep{bucholtz_community-centered_2021} and user centered design methods \citep{wallisch_overcoming_2019}, 
we involved target users early on in our explorations of the problem space to ensure we were 1) Solving a real issue faced by Tamil learners and 2) Providing a solution that learners would want to adopt. 
We wanted to understand what bottlenecks users were currently facing when searching for Tamil word definitions as well as existing paint points they were experiencing. 

\subsection{Prototype}
We built a medium fidelity prototype in Figma to test our design. 
We wanted to know whether a) users would be able to determine the meaning 
of a word when presented with the morphological breakdown and if 
b) users found the morphological breakdown useful and 
c) if users would use the tool if it existed.

\subsection{Participants}
To determine our design requirements, we interviewed 10 Tamil speakers 
who were asked to self describe Tamil language abilities in terms of oral fluency and literacy. 
Of our participants, 3 were native speakers, 6 were heritage speakers, and one was an L2 learner 
with no Tamil background. Importantly, oral fluency and literacy were not correlated for our speakers. 
In other words, speakers that were fluent in Tamil were not necessarily strong readers. 
As shown in (2), most of the participants had stronger oral abilities compared to reading and writing. 

\subsection{Exploratory Interview}
At the beginning of the session, participants were asked to self-describe their 
Tamil oral and literacy abilities. In addition, since most of the participants were heritage speakers, 
they were also asked demographic and language related questions to gain insight into their Tamil proficiency. 
As described by \citet{montrul_heritage_2016}, background information, 
such as level of education, language of instruction, domains of use of their languages, etc., 
is important to understanding a heritage speaker’s ability in their heritage language.
Although the answers to these questions were mostly used to corroborate the 
responder’s self description of their Tamil abilities and not used for any formal analyses, 
the questions themselves have been included in the appendix for future researchers. 
\\
Users were then asked to complete five tasks to test the usability of the prototype’s UI as well as input methods, 
which can also be found in the appendices.  

After completion of the tasks, users were asked the following questions:
\begin{itemize}
  \item What do you like/dislike about this tool?
  \item What features do you wish existed?
  \item How did each input method compare?
  \item How did the prototype compare to how you were originally looking up Tamil word definitions in the first task?
  \item Would you use this tool?
\end{itemize}

\subsection{Interview Results}
\textbf{Task 1:}\\
During the exploratory interviews, participants were asked what tools they were currently using when looking 
up the meaning of Tamil words. The most common answer given was Google Translate, 
however not all participants were using it just to get the translation of a word. 
For example, L, who had strong oral abilities but could not read Tamil easily, 
used Google Translate’s TTS feature to hear the word and would then be able to figure out what it meant. 
On the other hand, IP would use Google Translate to get a Tamil word’s romanization which they would 
then input into the Google search engine to try to find different contexts and uses of the word. 
For the participants who did rely on the translation provided by Google Translate, 
a couple Interestingly, none of our participants mentioned using generative AI tools,
however this could be due to the fact that most of our participants were over the age of twenty-five. 
As \citet{MAGOTRA201681} notes, younger individuals are more likely to 
adopt new technologies compared to older individuals.
\\
When asked if they would use this tool, seven participants responded with yes, two responded with no, and one with maybe.  
%insert image
SM, who is fluent with weak literacy skills, stated they would not use the tool. 
TB, who is a beginner heritage speaker that cannot read, stated the reason behind 
their “maybe” response was because they felt that their level of Tamil was 
too low to use the tool, but would use it if it included more conversational words. 
IP, who is fluent and native, felt the tool was unnecessary and that the 
interface made it hard to correlate meaning with the word parts.
\\
Conversely, the other participants felt that the tool was quite useful. 
In N’s case, they thought that learners would find the pedagogical 
breakdown helpful and wished the tool existed then they were attending 
Tamil language class during their adolescence. D, who self-described as 
conversationally fluent, thought that the pedagogical breakdown was 
not only helpful, but also corrects errors made by Google Translate. 
Lastly, R, who had only recently started learning Tamil thought that 
the definitions provided by the tool were clearer than Google Translate.

\section{User Requirements}
Based on the results of the exploratory interviews, we established the 
following user requirements to guide the design of the website. 
We also drew inspiration from the Kawennon:nis ´ verb conjugator for 
Kanyen'keha \citep{kazantseva_kawennonnis_2018} and its aim to 
assist users in acquiring morphology patterns rather than 
arbitrarily memorizing definitions of Tamil words. 
\begin{enumerate}
  \item Intuitively understand the morphological breakdown even as an L2 learner.
  \item Not overwhelm a beginner learner. 
  \item Is accessible on the phone. 
\end{enumerate}

\section{Systems Design}
\subsection{Romanized/Tamil Input}
% Users can input using the Tamil script or romanization (i.e., கேள்வி or kelvi). 
All of the dictionary and morphological data is stored in the Tamil script, 
so if a user inputs using romanization, their input is first converted into 
the Tamil script using a romanization-to-Tamil mapping written using 
G2P \citep{pine_gi2pi_2022} before being processed. 
Since there is no universal romanization standard followed by Tamil speakers, 
our mapping simply attempts to capture the most commonly used spelling patterns. 
After processing, the output is then returned to the user in both the 
Tamil script and romanization. The romanization presented here is derived 
from a Tamil-to-romanization mapping, again written using G2P. 
The romanization system used here is a combination of commonly used 
spelling conventions and IPA.

\subsection{Lemmatization}
Where possible, lemmatization of the Tamilinput is performed using 
Stanza \citep{qi2020stanza}. The output lemma is then checked 
against the dictionary data to see if it is an existing headword. 
If no valid dictionary entry is found, other strategies are 
employed to identify the lemma, such as suffix matching.

\subsection{Suffix Matching}
Where a valid lemma is not found by the lemmatizer, 
suffix matching is performed. First, a backwards search is 
performed on the input to see if there is a match with any
suffix included in the Gramble JSON file. If there are multiple matches, 
the longest one is selected. We create a new potential lemma by 
removing the text of the matched suffix from the input word. 
If the potential lemma is a valid lemma, we return this alongside the 
suffix as the output to the user. If the potential lemma is not a valid lemma, 
we check if the suffix that we found in the Gramble JSON has any 
material in its “add-back” field: this field contains phonological 
information that needs to be added to a stem in order to turn it into a lemma. 
If the “add-back” field is non-empty, we add this material to the potential lemma, 
confirm if this is now a valid lemma, and return it. If this strategy also fails, 
we return a “word not found” error to the user.

\subsection{Datasets}
We used two open source datasets for the Tamil dictionary component, specifically 
\citet{ylonen_wiktextract_2022} and \citet{vasuki_dharma_nodate}.

\subsection{Suffix Database Generation}

The suffix database queried by the system was created using Gramble. 
Documentation for Gramble can be found here: https://nrc-cnrc.github.io/gramble/; 
in this section we limit discussion to what tiers (that are then converted into JSON attributes) 
we created for suffix generation and why. All attributes are string type.
The primary attributes that every suffix entry in the suffix database have are 
\textit{text, gloss, display,} and \textit{add-back}. \textit{text} is a string representing a suffix 
as it appears in an input word, and is the attribute that the input is 
matched against as described in the system design walkthrough (section). 
% For instance, the input பூதங்கள் boothangaL ‘ghosts’ would match to the 
% suffix ங்கள் ngaL, as shown in Figure. 
\textit{gloss} is a layman’s definition of this suffix.
\textit{Display} is the version of the matched suffix that is shown to the user: 
% this is generally the same as what is in text, but differs in situations 
% where a suffix undergoes minor changes based on phonological environment. 
% For example, the plural suffix கள் gaL can also also appear as க்கள் kkaL or 
% ங்கள் ngaL based on the sound it is preceded by, but since கள் is its 
%  form, this is the form of the suffix that is shown (“displayed”) to the 
%  user in all cases, as shown in Figure. Finally, add-back is used as a 
%  secondary lemmatizer when the primary fails, as detailed above in section. 
%  In this instance, if the lemmatizer were unable to parse பூதங்கள், 
%  once the suffix match ங்கள் was identified and removed from the input, 
%  the character ம் (m) would then be “added back” to this new form to create 
%  the valid lemma பூதம் bootham (‘ghost’). 

\section{Results}
\subsection{Systems Analysis}
Due to limited resources and time (see Limitations chapter), 
we have not yet been able to perform a formal statistical evaluation. 
The minimal test set that was used to evaluate performance so far is 
composed of “edge-case” words, or words we predicted may be tricky to 
parse based on how our model is constructed. For instance, during early 
development, we noticed that the Stanford NLP lemmatizer often fails on 
class III verbs (as defined in Arden 1910), so we made sure to explicitly 
compensate for this during the suffix generation process using Gramble, 
and included several words from this category as part of our testing. 
Since we observed the remaining four verb classes \citep{arden_progressive_1910} being 
accurately lemmatized, fewer of these are represented in our test set 
as they were assumed to be higher performing. Similarly, there is a 
higher proportion of words with complex morphophonology and suffixal 
combinations present in our test set than is necessarily representative 
of the language, as we predicted these words to be the most challenging 
to cover. Thus, many “simple” words, for example words of the 
form [lemma + tense + person-number-gender], are not included in the test set, 
as once we confirmed that a sample few were working, we assumed others of the 
type would also be relatively high performing. Of course, formalized 
random testing is required to verify this. A high level overview of the 
tool’s performance on this test set is shown in table 1.
We primarily tested for presence of lemmas in our dictionary data (shown in table 2), 
and accurate lemmatization and suffix identification (shown in table 3). 
For inputs that produced a “word not found” error, we isolated which were 
failing because the lemma itself was not present in our dictionary data, 
and which were failing because our tool was not able to lemmatize the word, 
as shown in Table 4. This delineation will help us isolate which types of 
structures need to be added to our suffix database. 
A major category of words that are currently missing is lexical compounds. 
These are words that have two lemmas, which our model does not currently account for. 
% Some lexemes commonly used in compounds like போது po:thu (‘when’, colloquially) 
can be added to the database as suffixes, but Tamil also has a 
productive process of forming novel compounds, that will require a more nuanced approach.
For inputs that produced incorrect outputs, we differentiated between 
whether the error was in the lemma, the suffix, or both (table 5). 
The results from tables 4 and 5 in particular highlight that more 
work can be done in improving the suffix database, both to increase 
the accuracy and quantity of suffixes identified, as well as to assist 
in the lemmatization process. Thus this is where the focus of our future work lies.

\subsection{User Feedback}
We conducted a workshop for the launch of the MVP for Kelvi to 
gather user feedback. We also elicited feedback through word-of-mouth, 
and online through Reddit forums such as r/Tamil and r/LearningTamil. 
The feedback received was largely positive. Users requested voice input, 
romanization input, including example sentences in the output, and 
providing the definition of the input word in its entirety 
(not just of the individual morphemes). We were able to add the 
feature of romanization input following this feedback, and are 
evaluating the feasibility of incorporating the other requested features.

\section{Discussion}
Our current bottlenecks and proposed improvements to the software can be summarized as follows:
\begin{enumerate}
  \item We have limited root-word (i.e., dictionary) data, and results can be improved if we are able to acquire more data. 
  \item The lemmatizer that we use fails on some inputs; as discussed in (section), we are compensating for these gaps using Gramble, which we will continue to do. 
  \item Some morphemes are homophonous; for example, the suffix -um on a verb means “it will”, but on a noun means “and”. Currently we return both definitions, but if we are able to acquire reliable part-of-speech data, we can refine the output returned to the user. 
  \item Users are able to input using both the Tamil script and romanization, but the romanization input is significantly underperforming. This is because there is no romanization standard consistently followed by Tamil speakers, and so the mapping we currently employ cannot account for all possible spellings of a given Tamil word. The romanization practices of Tamil speakers are too irregular to be adequately captured by any one-to-many mapping, no matter how nuanced, and so this feature can only be improved if we are able to train a model on Tamil romanization input and/or implement an approximate search algorithm (\citealp{pine_zero-shot_2025}, \citealp{littell_waldayu_2017}).
  \item Since Kelvi was primarily designed to capture written forms, words in the spoken/informal register are underrepresented. Since these words vary widely across a range of dialects, collaboration with speakers across this span of variation would greatly enhance the coverage of our tool.
\end{enumerate}

\section{Conclusion}
Although Tamil has a relatively large speaker population, it is still considered a 
low-resource language which adds a layer of difficulty when developing digital tools 
and resources for Tamil that make use of data driven innovations in 
computational linguistics \citep{noauthor_tamil_2025} 
%(Tissera 2025). 
The system described in this paper serves as a starting point for 
creating pedagogical dictionaries designed for learners of 
polysynthetic languages and the morphological challenges they 
commonly face \citep{kazantseva_kawennonnis_2018}. 
Although there is room for improvement in terms of coverage and accuracy, 
we hope that Kelvi serves as an example of a grassroots initiative that, 
despite limitations in terms of language data and resources, 
is able to create a useful language learning tool by 
taking a user centered approach and targeting common pain points 
experienced by learners. The code and data used for the project 
is open source and will hopefully benefit other individuals and 
organizations looking to create language learning resources. 


% \section{Engines}

% To produce a PDF file, pdf\LaTeX{} is strongly recommended (over original \LaTeX{} plus dvips+ps2pdf or dvipdf).
% The style file \texttt{acl.sty} can also be used with
% lua\LaTeX{} and
% Xe\LaTeX{}, which are especially suitable for text in non-Latin scripts.
% The file \texttt{acl\_lualatex.tex} in this repository provides
% an example of how to use \texttt{acl.sty} with either
% lua\LaTeX{} or
% Xe\LaTeX{}.

% \section{Preamble}

% The first line of the file must be
% \begin{quote}
% \begin{verbatim}
% \documentclass[11pt]{article}
% \end{verbatim}
% \end{quote}

To load the style file in the review version:
\begin{quote}
\begin{verbatim}
\usepackage[review]{acl}
\end{verbatim}
\end{quote}
For the final version, omit the \verb|review| option:
\begin{quote}
\begin{verbatim}
\usepackage{acl}
\end{verbatim}
\end{quote}

To use Times Roman, put the following in the preamble:
\begin{quote}
\begin{verbatim}
\usepackage{times}
\end{verbatim}
\end{quote}
(Alternatives like txfonts or newtx are also acceptable.)

Please see the \LaTeX{} source of this document for comments on other packages that may be useful.

Set the title and author using \verb|\title| and \verb|\author|. Within the author list, format multiple authors using \verb|\and| and \verb|\And| and \verb|\AND|; please see the \LaTeX{} source for examples.

By default, the box containing the title and author names is set to the minimum of 5 cm. If you need more space, include the following in the preamble:
\begin{quote}
\begin{verbatim}
\setlength\titlebox{<dim>}
\end{verbatim}
\end{quote}
where \verb|<dim>| is replaced with a length. Do not set this length smaller than 5 cm.

\section{Document Body}

\subsection{Footnotes}

Footnotes are inserted with the \verb|\footnote| command.\footnote{This is a footnote.}

\subsection{Tables and figures}

See Table~\ref{tab:accents} for an example of a table and its caption.
\textbf{Do not override the default caption sizes.}

\begin{table}
  \centering
  \begin{tabular}{lc}
    \hline
    \textbf{Command} & \textbf{Output} \\
    \hline
    \verb|{\"a}|     & {\"a}           \\
    \verb|{\^e}|     & {\^e}           \\
    \verb|{\`i}|     & {\`i}           \\
    \verb|{\.I}|     & {\.I}           \\
    \verb|{\o}|      & {\o}            \\
    \verb|{\'u}|     & {\'u}           \\
    \verb|{\aa}|     & {\aa}           \\\hline
  \end{tabular}
  \begin{tabular}{lc}
    \hline
    \textbf{Command} & \textbf{Output} \\
    \hline
    \verb|{\c c}|    & {\c c}          \\
    \verb|{\u g}|    & {\u g}          \\
    \verb|{\l}|      & {\l}            \\
    \verb|{\~n}|     & {\~n}           \\
    \verb|{\H o}|    & {\H o}          \\
    \verb|{\v r}|    & {\v r}          \\
    \verb|{\ss}|     & {\ss}           \\
    \hline
  \end{tabular}
  \caption{Example commands for accented characters, to be used in, \emph{e.g.}, Bib\TeX{} entries.}
  \label{tab:accents}
\end{table}

As much as possible, fonts in figures should conform
to the document fonts. See Figure~\ref{fig:experiments} for an example of a figure and its caption.

Using the \verb|graphicx| package graphics files can be included within figure
environment at an appropriate point within the text.
The \verb|graphicx| package supports various optional arguments to control the
appearance of the figure.
You must include it explicitly in the \LaTeX{} preamble (after the
\verb|\documentclass| declaration and before \verb|\begin{document}|) using
\verb|\usepackage{graphicx}|.

\begin{figure}[t]
  \includegraphics[width=\columnwidth]{example-image-golden}
  \caption{A figure with a caption that runs for more than one line.
    Example image is usually available through the \texttt{mwe} package
    without even mentioning it in the preamble.}
  \label{fig:experiments}
\end{figure}

\begin{figure*}[t]
  \includegraphics[width=0.48\linewidth]{example-image-a} \hfill
  \includegraphics[width=0.48\linewidth]{example-image-b}
  \caption {A minimal working example to demonstrate how to place
    two images side-by-side.}
\end{figure*}

\subsection{Hyperlinks}

Users of older versions of \LaTeX{} may encounter the following error during compilation:
\begin{quote}
\verb|\pdfendlink| ended up in different nesting level than \verb|\pdfstartlink|.
\end{quote}
This happens when pdf\LaTeX{} is used and a citation splits across a page boundary. The best way to fix this is to upgrade \LaTeX{} to 2018-12-01 or later.

\subsection{Citations}

\begin{table*}
  \centering
  \begin{tabular}{lll}
    \hline
    \textbf{Output}           & \textbf{natbib command} & \textbf{ACL only command} \\
    \hline
    \citep{Gusfield:97}       & \verb|\citep|           &                           \\
    \citealp{Gusfield:97}     & \verb|\citealp|         &                           \\
    \citet{Gusfield:97}       & \verb|\citet|           &                           \\
    \citeyearpar{Gusfield:97} & \verb|\citeyearpar|     &                           \\
    \citeposs{Gusfield:97}    &                         & \verb|\citeposs|          \\
    \hline
  \end{tabular}
  \caption{\label{citation-guide}
    Citation commands supported by the style file.
    The style is based on the natbib package and supports all natbib citation commands.
    It also supports commands defined in previous ACL style files for compatibility.
  }
\end{table*}

Table~\ref{citation-guide} shows the syntax supported by the style files.
We encourage you to use the natbib styles.
You can use the command \verb|\citet| (cite in text) to get ``author (year)'' citations, like this citation to a paper by \citet{Gusfield:97}.
You can use the command \verb|\citep| (cite in parentheses) to get ``(author, year)'' citations \citep{Gusfield:97}.
You can use the command \verb|\citealp| (alternative cite without parentheses) to get ``author, year'' citations, which is useful for using citations within parentheses (e.g. \citealp{Gusfield:97}).

A possessive citation can be made with the command \verb|\citeposs|.
This is not a standard natbib command, so it is generally not compatible
with other style files.

\subsection{References}

\nocite{Ando2005,andrew2007scalable,rasooli-tetrault-2015}

% The \LaTeX{} and Bib\TeX{} style files provided roughly follow the American Psychological Association format.
% If your own bib file is named \texttt{custom.bib}, then placing the following before any appendices in your \LaTeX{} file will generate the references section for you:
% \begin{quote}
% \begin{verbatim}
% \bibliography{custom}
% \end{verbatim}
% \end{quote}

You can obtain the complete ACL Anthology as a Bib\TeX{} file from \url{https://aclweb.org/anthology/anthology.bib.gz}.
To include both the Anthology and your own .bib file, use the following instead of the above.
\begin{quote}
\begin{verbatim}
\bibliography{custom}
\end{verbatim}
\end{quote}

Please see Section~\ref{sec:bibtex} for information on preparing Bib\TeX{} files.

\subsection{Equations}

An example equation is shown below:
\begin{equation}
  \label{eq:example}
  A = \pi r^2
\end{equation}

Labels for equation numbers, sections, subsections, figures and tables
are all defined with the \verb|\label{label}| command and cross references
to them are made with the \verb|\ref{label}| command.

This an example cross-reference to Equation~\ref{eq:example}.

\subsection{Appendices}

Use \verb|\appendix| before any appendix section to switch the section numbering over to letters. See Appendix~\ref{sec:appendix} for an example.

\section{Bib\TeX{} Files}
\label{sec:bibtex}

Unicode cannot be used in Bib\TeX{} entries, and some ways of typing special characters can disrupt Bib\TeX's alphabetization. The recommended way of typing special characters is shown in Table~\ref{tab:accents}.

Please ensure that Bib\TeX{} records contain DOIs or URLs when possible, and for all the ACL materials that you reference.
Use the \verb|doi| field for DOIs and the \verb|url| field for URLs.
If a Bib\TeX{} entry has a URL or DOI field, the paper title in the references section will appear as a hyperlink to the paper, using the hyperref \LaTeX{} package.

\section*{Limitations}

This document does not cover the content requirements for ACL or any
other specific venue.  Check the author instructions for
information on
maximum page lengths, the required ``Limitations'' section,
and so on.

\section*{Acknowledgments}

This document has been adapted
by Steven Bethard, Ryan Cotterell and Rui Yan
from the instructions for earlier ACL and NAACL proceedings, including those for
ACL 2019 by Douwe Kiela and Ivan Vuli\'{c},
NAACL 2019 by Stephanie Lukin and Alla Roskovskaya,
ACL 2018 by Shay Cohen, Kevin Gimpel, and Wei Lu,
NAACL 2018 by Margaret Mitchell and Stephanie Lukin,
Bib\TeX{} suggestions for (NA)ACL 2017/2018 from Jason Eisner,
ACL 2017 by Dan Gildea and Min-Yen Kan,
NAACL 2017 by Margaret Mitchell,
ACL 2012 by Maggie Li and Michael White,
ACL 2010 by Jing-Shin Chang and Philipp Koehn,
ACL 2008 by Johanna D. Moore, Simone Teufel, James Allan, and Sadaoki Furui,
ACL 2005 by Hwee Tou Ng and Kemal Oflazer,
ACL 2002 by Eugene Charniak and Dekang Lin,
and earlier ACL and EACL formats written by several people, including
John Chen, Henry S. Thompson and Donald Walker.
Additional elements were taken from the formatting instructions of the \emph{International Joint Conference on Artificial Intelligence} and the \emph{Conference on Computer Vision and Pattern Recognition}.

% Bibliography entries for the entire Anthology, followed by custom entries
%\bibliography{custom,anthology-overleaf-1,anthology-overleaf-2}

% Custom bibliography entries only
\bibliography{custom}

\appendix

\section{Example Appendix}
\label{sec:appendix}

This is an appendix.

\end{document}
